\section*{Independent audit of the NS concentration and conic transport}

The audited objects are the explicit coordinates, differential operators,
complete leading-field vector residual, rank-four conic connection and
finite coefficient germ in \texttt{ns\_scaling\_bridge.tex}. This audit
does not prove the source manuscript's global Navier--Stokes theorem.
The complete audited source SHA is given in the accompanying JSON receipt;
the prior material extension and interface are separately hashed there.
No author files were edited.

\subsection*{Concentration derivatives and the full vector residual}

Retain $A=1/2+h$, $D=1/2-h$, $d=1-\xi^2$,
$\ell=1-2h\xi^2$, $Q>0$, $0<h<1/2$, $-1<\xi<1$ and $X>0$.
The map
\[
 (Q,\xi,X)\longmapsto
 (t,z,\rho)=(1-Q(1-\xi^2),Q^D\xi,\sqrt{2QX})
\]
has inverse differential equal to the three vector fields
\begin{align*}
 \partial_t&=-\ell^{-1}\partial_Q
      +\frac{D\xi}{Q\ell}\partial_\xi
      +\frac X{Q\ell}\partial_X,\\
 \partial_z&=\frac{2\xi Q^{1-D}}\ell\partial_Q
      +\frac d{Q^D\ell}\partial_\xi
      -\frac{2\xi X}{Q^D\ell}\partial_X,\\
 \partial_\rho&=\frac\rho Q\partial_X.
\end{align*}
Multiplication by the exact forward Jacobian proves these identities.
Its time--axial minor has determinant $-Q^D\ell$, or $Q^D\ell$
when $t$ is replaced by $\tau=1-t$. There is no sign change
missing in the first vector field. The independent checker verifies
that these three fields commute on $Q,\xi,X$.

Applying them to $Q^wf(X,\xi)$ yields exactly the author's
$T_w$, $Z_w$ and $R=2X\partial_X^2+2\partial_X$.
The second axial derivative is $Q^{w-2D}Z_{w-D}Z_wf$;
$w-2D=w-1+2h$. The second time derivative is
$Q^{w-2}T_{w-1}T_wf$. These identities were checked for an
arbitrary symbolic smooth profile, not a finite numerical sample.

The leading-field divergence is
\[
 \rho^{-1}\partial_\rho(\rho u_\rho^{(0)})
       +\partial_z u_z^{(0)}
   =Q^{-1}\bigl((V_0)_X+Z_{-A}U\bigr).
\]
Writing $I(X,\xi)=\int_0^XU(y,\xi)\,dy$ and integrating
$XU_X$ by parts gives exactly
\[
 V_0=\ell^{-1}(2\xi XU-2D\xi I-d I_\xi).
\]
The integration constant is zero at the regular axis: smooth
$V_0/X$ implies $V_0(0,\xi)=0$. Thus no unrecorded free radial
flux is omitted. At $X=0$ the formulas involving $1/X$ are
interpreted through the stated jointly smooth velocity-axis profiles;
smooth pressure there additionally requires a smooth value
$\Pi(0,\xi)$, which is not inferred from $\Pi_X$ alone. The displayed
coordinate and residual computation itself is on $X>0$.

For the material derivative, the radial contribution is
$Q^{w-1}V_0f_X$ and the axial contribution is
$Q^{w-A-D}UZ_wf=Q^{w-1}UZ_wf$.
Thus its exact operator is $Q^{w-1}M_w$, as stated.
The radial swirl and pressure-gradient terms both have weight
$-2A-1/2=-3/2-2h$ and cancel by
$\Pi_X=E^2/(2X)$. The remaining radial velocity is
$u_\rho^{(0)}=Q^{-1/2}V_0/\sqrt{2X}$, giving the two retained
radial weights $-3/2$ and $-3/2+2h$. The azimuthal frame term
is $Q^{-A-1}V_0E/(2X)$, and the radial and azimuthal
vector-Laplacian corrections are exactly the subtraction of
$(2X)^{-1}$ from $R$. The axial pressure weight is
$-2A-D=-A-1$. Substitution proves all three displayed vector
residuals, with no vanishing residual asserted. The independent
profile checker verifies their full expressions before and after
the one stated radial cancellation.

\subsection*{Complete conic derivations and both second derivatives}

Put $m=b^2/16-2s$, $n=m+4$, $\beta=b/4$ and work over
\[
 \mathcal B=\mathbb C[s,b,(mn)^{-1}],\qquad
 \mathcal E=\mathcal B[\eta,p]/(\eta^2-m,p^2-n).
\]
Since $2\eta$ and $2p$ are units, differentiating these relations
uniquely gives
\[
 \nabla_s\eta=-1/\eta,\quad\nabla_sp=-1/p,\qquad
 \nabla_b\eta=b/(16\eta),\quad\nabla_bp=b/(16p).
\]
These annihilate the defining relations, proving existence. Their
commutator kills the base, and differentiation of the two relations
then forces it to kill $\eta,p$, proving flatness. The independent
checker also verifies the Leibniz identity for every product of the
four basis elements $1,\eta,p,\eta p$.

In that quadratic basis the matrices are
\[
 H_s=\operatorname{diag}(0,-1/m,-1/n,-1/m-1/n),\qquad
 H_b=-bH_s/16.
\]
The change-of-coefficient matrix from $(1,r,p,rp)$ to this basis is
\[
 T=\begin{pmatrix}1&\beta&0&0\\0&1&0&0\\
       0&0&1&\beta\\0&0&0&1\end{pmatrix}.
\]
For a coefficient vector $f$, differentiating $Tf$ gives
\[
 \Gamma_i=T^{-1}H_iT+T^{-1}\partial_iT.
\]
This proves both original-basis matrices in the author text, including
the $1/4$ terms in $\Gamma_b$. The new $\Gamma_s$ equals exactly
the previous material extension's matrix called $\Gamma_b$ there;
that older subscript denotes the retained parameter $b$ in an
$s$-direction connection. It is not the new $b$-direction matrix.

Write $J_i=\partial_iH_i+H_i^2$. Differentiation using
$m_s=n_s=-2$, $m_b=n_b=b/8$ gives
\begin{align*}
 J_s&=\operatorname{diag}(0,-m^{-2},-n^{-2},
                   -m^{-2}-n^{-2}+2/(mn)),\\
 (J_b)_{11}&=(16m)^{-1}-b^2/(256m^2),\\
 (J_b)_{22}&=(16n)^{-1}-b^2/(256n^2),\\
 (J_b)_{33}&=(m^{-1}+n^{-1})/16
       -b^2(m^{-2}+n^{-2})/256+b^2/(128mn),
 \qquad(J_b)_{00}=0.
\end{align*}
Every off-diagonal entry is zero. This confirms both author squares.
Their exact original-basis constant terms are
\[
 Z_i=\partial_i\Gamma_i+\Gamma_i^2
     =T^{-1}(J_iT+2H_i\partial_iT+\partial_i^2T),
\]
so the derivative of the varying basis is retained also at second order.
For an independent complete expansion of the new axial square, define
\[
 u=-\frac b{64m}+\frac{b^3}{1024m^2},\qquad
 v=u+\frac b{32n}-\frac{b^3}{512mn}.
\]
Then its entire constant matrix in $(1,r,p,rp)$ is
\[
 Z_b=\begin{pmatrix}
 0&u&0&0\\0&(J_b)_{11}&0&0\\
 0&0&(J_b)_{22}&v\\0&0&0&(J_b)_{33}
 \end{pmatrix}.
\]
Direct differentiation and multiplication of the author's
$\Gamma_b$ verifies each entry. The full operator is
$\partial_b^2I_4+2\Gamma_b\partial_b+Z_b$.
The analogous $Z_s$ agrees entry for entry with the prior material
extension's complete square. The checker verifies both square-gauge
identities and
$\partial_s\Gamma_b-\partial_b\Gamma_s+
\Gamma_s\Gamma_b-\Gamma_b\Gamma_s=0$ exactly.

\subsection*{Domains, physical base change and finite coefficient germ}

The conic inverse formulas $q=(p+\eta)/2$,
$q^{-1}=(p-\eta)/2$ have product $(n-m)/4=1$.
Their sum and difference recover $p,\eta$, and
$s=b^2/32-\eta^2/2$ recovers the original base coordinate.
Only $m=0$ and $n=0$ have been removed from the finite conic
base. In Laurent coordinates this removes $q=\pm1,\pm i$;
$q=0$ is already absent from the Laurent presentation.

On $\Omega$ the base map $s=1-t$, $b=z$ is an explicit real
coordinate identification. Extending scalars to complex smooth
functions on this open base gives a rank-four algebra of sections.
The extended derivations act on those coefficients as
$\nabla_s=-\partial_t$, $\nabla_b=\partial_z$ and the radial
derivative kills the conic generators. The connection matrices
have no $\rho$ dependence. Flatness therefore proves that
\[
 -\nabla_s-\nu(\partial_\rho^2+\rho^{-1}\partial_\rho+\nabla_b^2)
 \quad\hbox{commutes with}\quad-\nabla_s^2/4
\]
for constant $\nu$. This is a smooth algebra-bundle identity. It
does not define differentiation on the earlier arbitrary topological
arithmetic quotient; no such unsupported descent is used.
The scalar coefficient inclusion $f\mapsto f1$ is the stated
derivative-intertwining algebra map. The physical leading material
derivative remains the separate $M_w$ formula, and is not identified
with the original material base derivation $-\partial_t+6c\partial_z$.

For $-4<m<0$ all four section values are
$\eta=\pm i\sqrt{-m}$ and $p=\pm\sqrt{m+4}$.
They are distinct and give $|q|=1$. Pointwise complex conjugation
sends the first sign to its opposite and fixes the second; hence
it sends $q$ to $q^{-1}$. No real nonzero $q$ can occur because
$(q-q^{-1})^2=m<0$. These statements concern section values over
the physical real base and retain all four sheets.

On the specified growth path, $b=0$, $s=\tau$, $m=-2\tau$.
For $0<\tau<2$ the sheet $q=e^{i\theta}$ has
$\theta=\arcsin\sqrt{\tau/2}\in(0,\pi/2)$.
The exact retained coefficient satisfies
\[
 C(e^{i\theta})=-4\sin(7\theta)\sin(3\theta)
 =\tau(-42+196\tau-280\tau^2+160\tau^3-32\tau^4).
\]
The finite Laurent sums for $[3]$ and $[7]$ give the complete
polynomial after $p^2=4-2\tau$; the independent checker retains
all factors. Both sine factors are positive for
$0<\theta<\pi/7$, proving the claimed strict negative interval
$0<\tau<2\sin^2(\pi/7)$. The endpoint $\tau=0$ remains outside
the localized connection; the scalar polynomial has its asserted
limit and order there without extending that connection.

For bounded $\varepsilon$, $e_0>0$ and $0<h<1/100$, the exact
product on the full stated scalar-germ class is
\[
 C(q)\tau^{-A}(e_0+\tau^{2h}\varepsilon(\tau))
       =\tau^D B(\tau)(e_0+\tau^{2h}\varepsilon(\tau)).
\]
Since $B(0)=-42$, $B(\tau)+42=O(\tau)$ and $0<2h<1$,
the error is $O(\tau^{D+2h})$, the leading coefficient is
$-42e_0$, and the product is negative for sufficiently small
positive $\tau$ while tending to zero. This proof requires neither
smoothness of the arbitrary bounded $\varepsilon$ nor a proof of
the manuscript's completed velocity claim. Finally the pressure
integral's sign is its stated infinity normalization: adding a
function of time preserves the gradient but can change that sign.

\paragraph{Audit conclusion.}
No mathematical discrepancy was found in the audited bridge.
The main independent checker verifies 374 exact conic, connection,
basis, square and coefficient identities. The delegated independent
profile checker verifies 34 exact inverse-coordinate, derivative,
divergence, material and full vector-residual identities. Their
combined scope supports the explicit formulas above, not an
independent Navier--Stokes blow-up theorem, a plethysm obstruction,
or a complexity separation.
